\section{GPT-3 与 ChatGPT：外部冲击如何改写公司判断}

认知智能给出了内部方向，但大模型浪潮不断从外部改写这条路线。GPT-3 和 ChatGPT 是两个不同层面的冲击：一个改变技术想象，一个改变用户和市场。

访谈章节把 GPT-3 和 ChatGPT 分开，这很重要。GPT-3 代表规模范式进入视野；ChatGPT 代表用户心智和产业节奏被突然点燃。张鹏回忆这些节点时，用的是“既焦虑又兴奋”的语气：焦虑来自外部浪潮太大，兴奋来自自己长期准备的方向终于被验证。

\begin{figure}[H]
\centering
\includegraphics[width=0.96\textwidth]{model-shockwaves.png}
\caption{三次模型冲击：GPT-3、ChatGPT、DeepSeek 及其后续范式变化。}
\end{figure}

\begin{knowledgebox}{读图：每次冲击都改变一个层面}
GPT-3 改变技术想象：规模可能带来通用能力。ChatGPT 改变用户心智：普通人知道大模型能做什么。DeepSeek 改变效率与开源叙事：模型能力不只靠砸资源，也可以靠工程和开放生态撬动。
\end{knowledgebox}

\subsection{为什么“机会留给有准备的人”}

讲完外部冲击后，本节回答一个常见误解：机会来了是不是只靠运气。张鹏的回答更接近日常积累，而不是神奇预测。

张鹏引用导师的话：机会像海上的木板，漂过来时你还得扑腾两下才能抓住。大模型浪潮很难精准预测，但可以长期准备：坚持正确方向，不断积累工程、研究和商业化能力，机会来时才有能力抓住。

\begin{importantbox}{准备不是预测}
预测是知道明天会发生什么；准备是即使不知道明天发生什么，也持续沿正确方向积累能力。大模型浪潮更像后者。
\end{importantbox}

\subsection{本章小结}

智谱的模型路线不是单纯跟风 GPT-3 或 ChatGPT，而是早期认知智能方向被外部技术范式反复校正。外部冲击越大，越考验组织是否有准备。

